\origpage{256--267}

\BellaSection{4. Extension d'un corps : construction de $\mathbb C$}

\medskip\noindent\textsc{Définition 7.50.} --- \emph{Un polynôme $P\in K[X]$ est dit scindé sur le corps $K$ s'il est constant ou bien si $P$ admet des zéros distincts $a_1,\ldots,a_r$ dans $K$ de multiplicités $k_1,\ldots,k_r$ avec $k_1+k_2+\cdots+k_r=n=$ degré de $P$.}

Cela revient à dire que $P$ admet une décomposition en produit de puissances de polynômes de degré 1, lorsqu'il est non constant.

\subsection*{Exemple}

\[
X^2-4\text{ est scindé sur }\mathbb Q,
\]
\[
X^2-2\text{ est scindé sur }\mathbb R\text{ mais pas sur }\mathbb Q.
\]

En effet, si $\sqrt2\in\mathbb Q$, en écrivant $\sqrt2=\dfrac pq$ irréductible, on aurait $2q^2=p^2$ donc 2 divise $p$, avec $p=2p_1$ on a $2q^2=4p_1^2\Rightarrow q^2=2p_1^2$ donc 2 divise $q$, mais alors $p$ et $q$ étant tous deux divisibles par 2, cela contredit $p/q$ irréductible.

Cet exemple, joint au fait que $\mathbb Q$ peut être considéré comme un sous-corps de $\mathbb R$, nous montre qu'un polynôme sera scindé ou non suivant le corps où sont les coefficients.

\medskip\noindent\textsc{Définition 7.51.} --- \emph{Un corps $K$ est dit algébriquement clos si tout polynôme non constant de $K[X]$ admet au moins un zéro dans $K$.}

\medskip\noindent\textsc{Théorème 7.52.} --- \emph{Un corps $K$ est algébriquement clos si et seulement si tout polynôme $P$ de $K[X]$ est scindé sur $K$, ou encore si et seulement si les seuls polynômes irréductibles sont de degré 1.}

Soit $K$ algébriquement clos, $P$ non constant dans $K[X]$, il admet un zéro $a_1$, donc (Théorème 7.38), $P$ est divisible par $X-a_1$, on a $P=(X-a_1)P_1$ avec $d^\circ P_1=(d^\circ P)-1$. Si $d^\circ P_1\ge1$, à son tour $P_1$ est divisible par $X-a_2$, on continue d'où une factorisation du type
\[
P=(X-a_1)\cdots(X-a_n)P_n,
\]
avec $n=d^\circ P$ et $P_n$ constante non nulle.

On regroupe les $a_i$ égaux d'où à la fois une factorisation de $P$ prouvant qu'il est scindé, mais aussi le fait que $d^\circ P>1\Rightarrow P$ non irréductible. Bien sûr, $P$ de degré 1 est irréductible.

Réciproquement : Si tout polynôme $P$ est scindé, $P$ non constant s'écrira sous la forme
\[
P=(\mathrm{cte})\prod_{i=1}^r(X-a_i)^{k_i}
\]
avec $k_1+\cdots+k_r=d^\circ P$ donc $P$ admettra des zéros.

En effet, la décomposition donnée en 7.35 dans un anneau principal nous dit que, si les seuls polynômes irréductibles sont de degré 1, on a une décomposition de tout polynôme non nul $P$ du type
\[
P=u\prod_{a\in K}(X-a)^{\alpha(a)},
\]
avec $u$ inversible dans $K[X]$, c'est-à-dire $u$ dans $K^*$, puisque pour des $a$ différents les $X-a$ ne sont pas associés; on a les $\alpha(a)$ entiers presque tous nuls. Là encore, si $n=\sum_{a\in K}\alpha(a)$ est non nul, $P$ admet au moins un zéro d'où la réciproque. \bookend

\medskip\noindent\textsc{Corollaire 7.53.} --- \emph{Si $K$ est algébriquement clos, $P$ non nul divise $Q$ si et seulement si tout zéro de $P$ est zéro de $Q$ avec une multiplicité supérieure ou égale.}

Puisque
\[
P=u\prod_{a\in K}(X-a)^{\alpha(a)}
\]
divisera
\[
Q=v\prod_{a\in K}(X-a)^{\beta(a)}
\]
si et seulement si, $\forall a\in K$, $\alpha(a)\le\beta(a)$. \bookend

Le problème qui se pose, c'est que les corps ne sont pas tous algébriquement clos. Ainsi, $\mathbb Q$ ne l'est pas, ($X^2-2$ non scindé), $\mathbb R$ non plus, ($X^2+1$ non scindé). Comment faire?

\medskip\noindent\textsc{Définition 7.54.} --- \emph{On appelle extension du corps commutatif $K$, tout corps $L$ contenant $K$ comme sous-corps.}

Ainsi, en identifiant $\mathbb Q$ avec son image injective dans $\mathbb R$ (voir chapitre 6 de Topologie) on peut dire que $\mathbb R$ est un sur-corps de $\mathbb Q$.

\medskip\noindent\textsc{Définition 7.55.} --- \emph{Soit $L$ une extension d'un corps $K$ et $a\in L$. On appelle extension simple de $K$ par $a$ dans $L$ le plus petit sous-corps de $L$ contenant $K$ et $a$.}

On la note $K(a)$ et $K(a)$ est visiblement l'intersection de tous les sous-corps de $L$ contenant $K$ et $a$. (Il y en a, ne serait ce que $L$).

On a $K(a)=K\Longleftrightarrow a\in K$.

Soit alors $\theta:K[X]\longrightarrow K(a)$ l'application qui au polynôme
\[
P(X)=\sum_{i=0}^{d^\circ P}a_iX^i
\]
associe
\[
\theta(P)=\sum_{i=0}^{d^\circ P}a_i a^i.
\]
Comme les $a_i$ et $a$ sont dans le corps $K(a)$, $\theta(P)\in K(a)$.

Il est immédiat de vérifier que $\theta$ est un morphisme d'anneau de $K[X]$ dans l'anneau $K(a)$. Son noyau $\operatorname{Ker}\theta$ est un idéal de $K[X]$, anneau principal. Donc l'image $\theta(K[X])$ est un anneau isomorphe à l'anneau quotient $K[X]/\operatorname{Ker}\theta$. Il y a alors deux cas de figure.

\subsection*{7.56.}

Si $\operatorname{Ker}\theta=\{0\}$, (donc pour tout polynôme $P$ non nul de $K[X]$ on a $P(a)\ne0$ : l'élément $a$ est alors dit transcendant sur $K$), $K[X]/\operatorname{Ker}\theta\simeq\{P(a);P\in K[X]\}$ est un anneau, (sous-anneau de $K(a)$) intègre, $K[X]$ étant intègre et $\theta$ bijective de $K[X]$ sur cet anneau que l'on notera $K[a]$. Mais alors $K[a]$ admet dans le sur-corps $L$ son corps des fractions (Théorème 5.34) à un isomorphisme près on peut donc dire que \BellaCorrectionNote{$K(a)$}{$L(a)$}{Comme $a\in L$, on aurait $L(a)=L$; le corps des fractions de $K[a]$ est l'extension simple $K(a)$.} est le corps des fractions de l'anneau $K[a]$.

\subsection*{7.57.}

\BellaCorrectionNote{Si}{5.57.}{La numérotation imprimée rompt la suite 7.56, 7.58; il s’agit du numéro 7.57.} $\operatorname{Ker}\theta\ne\{0\}$, soit $P$ le polynôme unitaire engendrant $\operatorname{Ker}\theta$. L'élément $a$ est alors dit algébrique sur $K$, l'extension simple $K(a)$ est dite algébrique sur $K$ et $P$ est le polynôme minimal de $a$ sur $K$.

\medskip\noindent\textsc{Théorème 7.58.} --- \emph{Soit $a$ algébrique sur un corps $K$. Son polynôme minimal est irréductible sur $K$. Il est de degré 2 au moins si $a\notin K$ et $K(a)$ est alors un espace vectoriel de dimension $n$ sur $K$, $n$ degré du polynôme minimal.}

Pour l'instant, on suppose connu un sur-corps $L$ et tout dépend encore de $L$.

Soit donc $a$ algébrique sur un corps $K$ et $L$ un sur-corps dans lequel on considère l'extension simple $K(a)$.

Si $P$ engendre l'idéal $\operatorname{Ker}\theta$, et si $P$ est non irréductible, avec $P(X)=P_1(X)P_2(X)$, puisque $\theta(P)=P(a)=0$, c'est que dans le corps $L$, on a $P_1(a)$ ou $P_2(a)=0$, donc $P_1$ (ou $P_2$) $\in\operatorname{Ker}\theta$, on aurait $P_1$ (ou $P_2$) multiple de $P$ avec un degré strictement inférieur : c'est exclu. D'où $P$ irréductible.

Si $d^\circ P=1$, si $P(X)=\alpha X+\beta$ avec $\alpha\ne0$, comme $\alpha a+\beta=0$ on aurait $a=-\alpha^{-1}\beta\in K$, et si $a\in K$, $X-a$ est un polynôme de $K[X]$ qui annule $a$ donc $\operatorname{Ker}\theta$ est engendré par $X-a$ : le polynôme minimal est bien de degré 1 si et seulement si $a\in K$.

Enfin $P$ étant irréductible, l'idéal qu'il engendre est maximal (Théorème 7.24) donc l'anneau quotient $K[X]/\operatorname{Ker}\theta$ est un corps (Théorème 5.43). Comme ce quotient est isomorphe à \BellaCorrectionNote{$\theta(K[X])=K[a]$}{$\theta(K)=K[a]$}{L'image de $K$ par $\theta$ ne contient que les constantes; l'anneau $K[a]$ est l'image de tout $K[X]$.} c'est que $K[a]$ est un corps, contenant $K$, et $a$, donc contenant $K(a)$, plus petit corps de $L$ contenant $K$ et $a$. Comme on a vu que $\theta$ est à valeurs dans $K(a)$, dans ce cas on a $K(a)=K[a]$.

De plus, pour tout $Q\in K[X]$, la division euclidienne de $Q$ par $P$ polynôme minimal de $a$ donne l'identité
\[
Q(X)=P(X)S(X)+R(X)
\]
avec $d^\circ R<d^\circ P$. Donc $Q(a)=R(a)$ puisque $P(a)=0$.

Il en résulte que
\[
\theta(Q)=R(a)\in\left\{\sum_{i=0}^{n-1}a_i a^i;\ a_i\in K\right\}
\]
avec $n=d^\circ P$.

Si on avait $R(a)=0$ avec $d^\circ R<n$, c'est que $R\in\operatorname{Ker}\theta$ idéal engendré par $P$ donc $R=0\cdot P$. Il en résulte que, dans la structure de $L$, espace vectoriel sur $K$, les éléments $1,a,a^2,\ldots,a^{n-1}$ sont indépendants et alors $K[a]=K(a)$ est en fait un espace vectoriel de dimension $n$ sur $K$. \bookend

Il reste une question à régler. On est parti de $L$ sur-corps de $K$, on a considéré $a\in L$ et l'extension simple $K(a)$ qui est transcendante ou algébrique.

Si on part de $L'$ autre sur-corps de $K$ contenant aussi $a$, on aura dans $L'$ l'extension simple notée $K'(a)$, (bien que l'on parte du même corps $K$).

A-t'on encore une extension algébrique si elle l'est sur $K$? Et si oui, quel est le degré du polynôme minimal?

\subsection*{7.59.}

On peut déjà remarquer que, si $a$ est \BellaCorrection{transcendant}{tanscendant} sur $K$ dans $L$ le corps $K(a)$ est espace vectoriel de dimension infinie sur $K$, car pour tout $n$, les éléments $1,a,a^2,\ldots,a^n$ sont libres sur $K$, une combinaison linéaire nulle non triviale fournissant un polynôme non nul dans l'idéal $\operatorname{Ker}\theta$.

\subsection*{7.60.}

Mais alors, il en résulte que si $a$ est algébrique sur $K$ dans une extension $L$ de $K$, il l'est dans toute extension $L'$ de $K$ le contenant, et il en est de même pour $a$ transcendant. De plus le degré du polynôme minimal sera le même.

En effet ce polynôme minimal, lui a ses coefficients dans $K$ donc si on notait

$P$ le polynôme minimal pour l'extension $L$ et

$P'$ le polynôme minimal pour l'extension $L'$

on aurait $P(a)=P'(a)=0$ donc $P'\in\operatorname{Ker}\theta$ donc $P$ divise $P'$, et $P\in\operatorname{Ker}\theta'$ d'où $P'$ divise $P$ et comme $P$ et $P'$ sont unitaires, on a $P=P'$.

\subsection*{7.61.}

Le degré $n$ de ce polynôme minimal est alors le degré de l'élément $a$ sur le corps $K$.

Nous venons donc de justifier que, si $L$ est un sur-corps de $K$, et si $a\in L\setminus K$, le plus petit sous-corps $K(a)$ de $L$ contenant $K$ et $a$ est de deux types possibles :

- soit $K(a)$ est espace vectoriel de dimension finie $n$ sur $K$ et alors $a$ est un zéro dans $K(a)$ d'un polynôme $P$ irréductible de $K[X]$ de degré $n\ge2$,

- soit $K(a)$ est le corps des fractions de $K[a]$ et il sera de dimension infinie en tant qu'espace vectoriel sur $K$ et aucun polynôme $P$ de $K[X]$ n'admet $a$ pour zéro.

On peut considérer le problème dans l'autre sens et, partant de $P\in K[X]$, $P$ irréductible de degré 2 au moins, \BellaCorrectionNote{constituer un (ou des) sur-corps de $K$}{constitue un (ou des) sur-corps de $L$}{Le corps que l'on cherche doit être une extension du corps de départ $K$; $L$ n'est pas encore défini à cet endroit.} dans lequel $P$ admettra au moins un zéro.

\subsection*{7.62.}

La méthode de l'adjonction symbolique due à Cauchy permet de résoudre cette question.

\medskip\noindent\textsc{Théorème 7.63.} --- \emph{Soit un corps $K$ et un polynôme $P\in K[X]$ irréductible de degré 2 au moins. Il existe au moins une extension algébrique simple $L$ de $K$ contenant un zéro $a$ de $P$. Si $L_1$ est une autre extension algébrique simple de $K$ contenant un zéro de $P$, alors $L$ et $L_1$ sont isomorphes.}

Le polynôme $P$ est irréductible donc l'idéal principal $(P)$ qu'il engendre est maximal (Théorème 7.24), l'anneau quotient
\[
L=K[X]/(P)
\]
est donc un corps (Théorème 5.43).

Soit $\varphi$ l'application de $K$ dans $L$ définie par :
\[
\forall k\in K,\qquad \varphi(k)=\text{la classe du polynôme constant }k\text{ dans }L.
\]
On sait que $\varphi$ est un morphisme d'anneaux. Il est injectif car
\[
\varphi(k)=\varphi(k')\Longleftrightarrow k-k'\text{ multiple de }P,
\]
\[
\Longleftrightarrow k-k'=PQ,
\]
avec $k-k'$ de degré $\le0$ et si $Q\ne0$, $d^\circ(PQ)\ge d^\circ P\ge2$. C'est donc que $Q=0$ d'où $k=k'$.

Si on identifie $K$ et son image $\varphi(K)$ dans $L$ on peut dire que $L$ devient un sur-corps de $K$. On retrouve le type de situation qui conduit à dire que $\mathbb N$ est contenu dans $\mathbb Z$ lui-même « contenu » dans $\mathbb Q$.

Soit alors $a$ la classe du polynôme $X$ dans le corps $L$. Si le polynôme $P$ s'écrit
\[
P(X)=\alpha_0+\alpha_1X+\cdots+\alpha_nX^n,
\]
(avec les $\alpha_n\in K$) la structure quotient implique que 0, classe de $P$ dans $L$ vérifie
\[
0=\varphi(\alpha_0)+\varphi(\alpha_1)a+\cdots+\varphi(\alpha_n)a^n
\]
ou encore, comme les $\varphi(\alpha_j)$ sont identifiés aux éléments $\alpha_j$ de $K$, que
\[
0=\alpha_0+\alpha_1a+\cdots+\alpha_na^n :\qquad\text{dans }L,\text{ le polynôme }P\text{ admet }a\text{ pour zéro.}
\]

Comme pour $A\in K[X]$, la division euclidienne de $A$ par $P$ conduit à une relation du type
\[
A(X)=P(X)S(X)+R(X)\qquad\text{avec }d^\circ R<d^\circ P,
\]
si on a
\[
R(X)=r_0+r_1X+\cdots+r_kX^k,\qquad(k<n)
\]
il vient
\[
\operatorname{classe}(A)=r_0+r_1a+\cdots+r_ka^k,
\]
d'où $L$ est l'ensemble des éléments du type
\[
r_0+r_1a+\cdots+r_{n-1}a^{n-1}
\]
avec $r_j\in K$ : $L=K(a)$ est extension algébrique simple de $K$ et contient un zéro $a$ de $P$.

De même que pour le théorème 7.58, on justifierait que si $L_1=K(b)$ est une autre extension algébrique simple contenant un zéro $b$ de $P$ irréductible, on a $L$ et $L_1$ isomorphes par la bijection $\psi$ qui à l'élément générique
\[
r_0+r_1a+\cdots+r_{n-1}a^{n-1}
\]
de $L$ associe
\[
r_0+r_1b+\cdots+r_{n-1}b^{n-1}
\]
dans $L_1=K(b)$. \bookend

\medskip\noindent\textsc{Définition 7.64.} --- \emph{Un corps $L$ construit comme au théorème 7.63 s'appelle un corps de rupture du polynôme $P$.}

Soit alors $P$ irréductible sur $K$, de degré $\ge2$. On construit un corps de rupture $L_1$ de $P$; dans ce corps $P$ admet au moins un zéro $a$ donc $X-a$ divise $P\in L_1[X]$. En fait $P$ n'est plus irréductible dans $L_1[X]$. On va justifier, par récurrence sur le degré de $P$, le fait qu'il existe un sur-corps de $K$ dans lequel $P$ sera scindé.

\medskip\noindent\textsc{Définition 7.65.} --- \emph{Soit $P\in K[X]$. On appelle corps de décomposition, ou corps des racines, de $P$ tout sur-corps $L$ de $K$ dans lequel $P$ est scindé.}

\medskip\noindent\textsc{Théorème 7.66.} --- \emph{Tout polynôme $P$ de $K[X]$ admet des corps de décomposition.}

On procède par récurrence sur $d^\circ P$.

Si $d^\circ P=0$, $P$ est scindé sur $K$ qui convient. Il en est de même si $d^\circ P=1$.

Supposons le résultat acquis pour des polynômes de degré $n-1$ au plus.

Soit $P$ de degré $n$, et $Q$ un diviseur irréductible dans $K[X]$ de $P$, (éventuellement $Q=P$ si $P$ est irréductible). On a vu que $Q$ admet un corps de rupture $L_1$, dans lequel existe $a$ racine de $Q$ avec $a\notin K$ si $d^\circ Q\ge2$. Donc il existe $k$ entier $\ge1$ tel que dans $L_1[X]$, $(X-a)^k$ divise $Q$ donc $P$. En écrivant
\[
P=(X-a)^kP_1
\]
on a $P_1\in L_1[X]$, $d^\circ P_1<n$ donc $P_1$ admet un corps de décomposition $L_2$, sur-corps de $L_1$ donc aussi de $K$.

Dans $L_2[X]$, $P_1$ est scindé, $a\in L_1\subset L_2\Rightarrow$ dans $L_2[X]$,
\[
P(X)=(X-a)^kP_1(X)
\]
est scindé. \bookend

\medskip\noindent\textsc{Remarque 7.67.} --- \emph{Le fait de savoir que tout polynôme de $K[X]$ admet des corps de décomposition permet de faire certains raisonnements du type suivant.}

Soit $A\in M_n(K)$, matrice carrée d'ordre $n$ sur $K$, de polynôme caractéristique $P$, et $L$ un corps de décomposition de $P$. La matrice $A$ considérée dans $M_n(L)$, (possible car $K$ se plonge dans $L$) ayant son polynôme caractéristique scindé est trigonalisable, donc en travaillant sur la forme triangulaire, on justifie que si $(\lambda_1,\ldots,\lambda_n)$ est le spectre de $A$ dans $L$, pour tout polynôme $Q$, la matrice $Q(A)$ aura $Q(\lambda_1),\ldots,Q(\lambda_n)$ pour spectre. Or la matrice $Q(A)$ se calcule dans $M_n(K)$ si $Q$ est à coefficients dans $K$, donc son spectre ne dépend pas du voyage fait dans le sur-corps, et \BellaCorrectionNote{le spectre de $Q(A)$ sera, dans $L$, l'ensemble des $Q(\lambda_i)$, qui ne sont pas nécessairement dans $K$.}{le spectre de $Q(A)$ sera, dans $K$, l'ensemble des $Q(\lambda_i)\in K$.}{Une matrice à coefficients dans $K$ peut avoir des valeurs propres hors de $K$; par exemple une rotation réelle peut avoir des valeurs propres complexes. Dans un corps de décomposition $L$, le spectre est bien formé des $Q(\lambda_i)$.} (Attention : les $\lambda_i$ ne sont pas forcément toutes dans $K$).

\medskip\noindent\textsc{Définition 7.68.} --- \emph{Une extension $L$ de $K$ est dite algébrique sur $K$ si tout élément de $L$ est algébrique sur $K$.}

Par exemple $\mathbb R$ n'est pas extension algébrique de $\mathbb Q$. C'est ainsi que Hermite a démontré en 1873 que $e$ est transcendant, Lindemann a établi la transcendance de $\pi$ en 1882, et Liouville celle des nombres du \BellaCorrection{type}{types}
\[
\sum_{k=0}^{\infty}\frac1{10^{a_k}},\qquad
0<a_0<a_1<\cdots<a_k<\cdots\quad\text{et}\quad
\lim_{k\to+\infty}\frac{a_{k+1}}{a_k}=+\infty.
\]
Justifier la transcendance d'un réel n'est pas chose facile!

\medskip\noindent\textsc{Définition 7.69.} --- \emph{On appelle clôture algébrique $L$ d'un corps $K$, toute \BellaCorrectionNote{extension algébrique $L$ de $K$ qui soit un corps algébriquement clos}{extension $L$ de $K$ qui soit un corps algébriquement clos}{Sans l'hypothèse que l'extension soit algébrique, un corps algébriquement clos contenant $K$ ne serait pas nécessairement une clôture algébrique de $K$; par exemple $\mathbb C/\mathbb Q$ n'est pas algébrique.}.}

Par utilisation de l'axiome de Zorn, on justifie alors le résultat important suivant :

\medskip\noindent\textsc{Théorème 7.70.} (De Steinitz). --- \emph{Tout corps $K$ admet une clôture algébrique.}

Résultat que nous admettrons.

Ce résultat qui affirme l'existence de la clôture algébrique ne la caractérise pas. Dans le cas du corps $\mathbb R$ au départ, il se trouve que la simple construction de $\mathbb C$ corps de rupture de $X^2+1$ nous fournit une clôture algébrique de $\mathbb R$. Ce résultat découvert par d'Alembert s'appelle encore parfois théorème fondamental de l'algèbre bien que ses justifications fassent intervenir de l'analyse.

Construisons d'abord $\mathbb C$. Le polynôme $X^2+1$ est irréductible sur $\mathbb R$, car tout $\lambda$ réel a un carré $\lambda^2\ge0$, donc différent de $-1$.

Soit $\mathbb C$ un corps de rupture de $X^2+1$ sur $\mathbb R$, $\mathbb C$ \BellaCorrection{est}{et} une extension algébrique de degré 2 sur $\mathbb R$ et si $i$ désigne une racine de $X^2+1$ dans $\mathbb C$, les éléments de $\mathbb C$ s'écrivent $a+ib$ avec $a$ et $b$ réels. De plus $\mathbb C$ est espace vectoriel de dimension 2 sur $\mathbb R$, avec l'addition définie par
\[
a+ib+(a'+ib')=(a+a')+i(b+b').
\]

Pour le produit, on a
\[
(a+ib)(a'+ib')=aa'+i(ba'+ab')+i^2bb',
\]
or $i^2=-1$, donc le produit s'écrit encore
\[
(a+ib)(a'+ib')=(aa'-bb')+i(ab'+ba').
\]

On retrouve les règles usuelles, que l'on pouvait prendre pour définir une structure sur $\mathbb R^2$, et vérifier ainsi que l'on a un corps alors qu'ici on sait que l'on a un corps.

On suppose donc acquises toutes les propriétés usuelles sur $\mathbb C$.

\medskip\noindent\textsc{Théorème 7.71.} (de d'Alembert). --- \emph{$\mathbb C$ est un corps algébriquement clos.}

Il en existe une justification par l'analyse, faisant intervenir les propriétés des séries entières de variable complexe, et la formule de Green Riemann, je la donnerai après l'étude des séries entières.

Il en existe une justification dite « algébrique », due à Gauss, d'où le nom de théorème de d'Alembert-Gauss qu'on lui donne parfois. J'ai mis le algébrique entre guillemets car elle utilise de l'analyse, à savoir le

\medskip\noindent\textsc{Lemme 7.72.} --- \emph{Soit $P\in\mathbb R[X]$ un polynôme de degré impair, il s'annule sur $\mathbb R$.}

En effet $P$ étant de degré impair a des limites infinies de signes opposés en $+\infty$ et $-\infty$ donc il s'annule sur $\mathbb R$, (théorème des valeurs intermédiaires pour une fonction continue.) \bookend

Pour justifier le théorème de d'Alembert, à ce stade il nous faut admettre des résultats sur les polynômes de plusieurs variables, symétriques.

Soit $Q(X_1,\ldots,X_n)\in K[X_1,\ldots,X_n]$ un polynôme de plusieurs variables. (Je sais, je n'ai pas dit ce que c'est, tant pis, allez voir ailleurs). Il est dit symétrique si, $\forall\sigma\in S_n$, $Q(X_{\sigma(1)},\ldots,X_{\sigma(n)})=Q(X_1,\ldots,X_n)$. Et dans ce cas il existe un autre polynôme $R\in K[X_1,\ldots,X_n]$, tel que
\[
Q(X_1,\ldots,X_n)=R(\sigma_1,\sigma_2,\ldots,\sigma_n)
\]
avec les $\sigma_j$ fonctions élémentaires symétriques des $X_j$, à savoir :
\[
\sigma_1=X_1+X_2+\cdots+X_n,
\]
\[
\sigma_2=X_1X_2+X_1X_3+\cdots+X_1X_n+X_2X_3+\cdots+X_nX_{n-1}
\]
$=$ somme des produits 2 à 2,
\[
\sigma_3=\text{somme des produits 3 à 3},
\]
et pour finir
\[
\sigma_n=X_1\cdots X_n,\quad\text{produit des }n\text{ variables.}
\]

Il y a donc en fait tout un travail à fournir pour avoir le théorème de d'Alembert. Personnellement, j'aime mieux la justification par l'analyse...

Soit alors $Q\in K[X]$ un polynôme scindé. Si on écrit à la fois
\[
Q(X)=a_0X^n+a_1X^{n-1}+\cdots+a_n,\qquad a_0\ne0
\]
et
\[
Q(X)=a_0\prod_{i=1}^n(X-\lambda_i)
\]
en mettant en évidence les racines de $Q$ distinctes ou non, on aura
\[
\begin{aligned}
Q(X)=a_0[&X^n-(\lambda_1+\cdots+\lambda_n)X^{n-1}\\
&+(\sum\text{ produits 2 à 2 des }\lambda_i)X^{n-2}\\
&-(\sum\text{ produits 3 à 3 des }\lambda_i)X^{n-3}\\
&+\cdots+(-1)^n\lambda_1\cdots\lambda_n].
\end{aligned}
\]
d'où les relations dites symétriques entre les racines de $Q$ :
\[
\tag{7.73}
\sigma_1=-\frac{a_1}{a_0},\quad
\sigma_2=(-1)^2\frac{a_2}{a_0},\quad
\sigma_3=(-1)^3\frac{a_3}{a_0},\quad\ldots,\quad
\sigma_n=(-1)^n\frac{a_n}{a_0}.
\]

Il faut remarquer que, si sur le corps $K$ le polynôme $Q$ n'est pas scindé on peut introduire $\overline K$ corps de décomposition \BellaCorrectionNote{de $Q$}{de $K$}{Le corps de décomposition est ici choisi pour le polynôme $Q$, afin de pouvoir le factoriser.} , factoriser $Q$ dans $\overline K[X]$ pour justifier les relations précédentes et obtenir ainsi le fait que même si les $\lambda_j$ ne sont pas dans $K$, leurs fonctions symétriques y seront car les $(-1)^k\dfrac{a_k}{a_0}$ sont dans $K$.

Passons alors à la justification de d'Alembert-Gauss. Soit un polynôme
\[
P(X)=a_0X^n+a_1X^{n-1}+\cdots+a_n\in\mathbb C[X],
\]
on introduit
\[
\overline P(X)=\overline a_0X^n+\overline a_1X^{n-1}+\cdots+\overline a_n,
\]
(avec $\overline a_k=\alpha_k-i\beta_k$ si $a_k=\alpha_k+i\beta_k$, $\alpha_k$ et $\beta_k$ étant réels : c'est le conjugué de $a_k$).

Le polynôme $Q=P\overline P$ est à coefficients réels, car c'est un polynôme de degré $2n$, et pour $j\in[0,2n]$, le coefficient $b_{2n-j}$ de $X^j$ dans $Q$ est
\[
b_{2n-j}=\sum_{\substack{k,l\in[0,n]\\k+l=2n-j}}a_k\overline a_l,
\]
son conjugué
\[
\overline b_{2n-j}=\sum_{\substack{k,l\in[0,n]\\k+l=2n-j}}\overline a_k a_l,
\]
c'est donc $b_{2n-j}$ vu le rôle symétrique des indices.

Le degré de $Q$ peut s'écrire $2n=2^mn'$ avec en fait $n'$ impair et $m$ entier $\ge1$. Si on prouve que $Q$ a une racine complexe $z$ la relation
\[
Q(z)=P(z)\overline P(z)=0
\]
donnera $P(z)=0$ ou $\overline P(z)=0$, mais en conjugant, dans le 2e cas, on a $P(\overline z)=0$ : dans les 2 cas $P$ aura un zéro dans $\mathbb C$. Il nous reste donc à justifier l'existence dans $\mathbb C$ d'un zéro de $Q$.

\emph{On va justifier ceci par récurrence sur $m$.}

Soit $L$ un corps de décomposition sur $\mathbb C$, de $Q$ et $z_1,\ldots,z_{2n}$ les zéros, (distincts ou non) de $Q$ dans $L$. On a, dans $L[X]$, la factorisation
\[
Q(X)=a_0\overline a_0\prod_{j=1}^{2n}(X-z_j).
\]

Supposons donnés $N$ nombres réels distincts $c_1,\ldots,c_N$. Pour chaque indice $k\le N$, soit le polynôme
\[
Q_k(X)=\prod_{1\le i<j\le2n}\left(X-(z_i z_j+c_k(z_i+z_j))\right).
\]

C'est un polynôme symétrique en $z_1,z_2,\ldots,z_{2n}$ racines de $Q$, à coefficients réels, (1 est coefficient de $X$, $-1$ celui de $z_iz_j$ et $-c_k$ celui de $z_i$ et de $z_j$), donc d'après le résultat admis, $Q_k$ est un polynôme à coefficients réels par rapport aux fonctions symétriques élémentaires des $z_j$, donc finalement des $(-1)^r\dfrac{b_r}{b_0}$, avec les $b_r$ coefficients réels de $Q$ : on obtient $Q_k$ polynôme en $X$ à coefficients réels. Quel est son degré : c'est le nombre de parties à 2 éléments dans $\{1,\ldots,2n\}$, donc
\[
C_{2n}^2=\frac{2n(2n-1)}2
\]
soit
\[
n(2n-1)=2^{m-1}\cdot n'(2n-1)
\]
avec $n'(2n-1)$ impair.

Supposons d'abord $m=1$. Alors chaque $Q_k$, de degré impair admet un zéro réel, donc complexe.

Soit $t_k$ ce zéro : il existe $i_k$ et $j_k$ avec \BellaCorrectionNote{$1\le i_k<j_k\le2n$}{$1<i_k<j_k\le2n$}{Le produit définissant $Q_k$ porte sur toutes les paires $1\le i<j\le2n$; l’indice $i_k=1$ n’est donc pas exclu.} et
\[
t_k=z_{i_k}z_{j_k}+c_k(z_{i_k}+z_{j_k}).
\]

Si on suppose $N>n(2n-1)$, il existe forcément deux indices $k$ et $k'$ distincts associés au même couple $(i_k,j_k)$ avec $1\le i_k<j_k\le2n$. On a donc
\[
t_k=z_{i_k}z_{j_k}+c_k(z_{i_k}+z_{j_k})
\]
et
\[
t_{k'}=z_{i_k}z_{j_k}+c_{k'}(z_{i_k}+z_{j_k})
\]
réels donc complexes.

On a un système linéaire permettant de calculer $a=z_{i_k}+z_{j_k}$ et $b=z_{i_k}z_{j_k}$. Les $c_k$ étant distincts, $c_k-c_{k'}\ne0$ et on trouve
\[
a=\frac{t_k-t_{k'}}{c_k-c_{k'}},
\qquad
b=\frac{c_{k'}t_k-c_kt_{k'}}{c_{k'}-c_k},
\]
relations qui prouvent que $a$ et $b$ sont complexes et non dans $L\setminus\mathbb C$.

Mais alors $z_{i_k}$ et $z_{j_k}$ sont solutions de
\[
(X-z_{i_k})(X-z_{j_k})=X^2-aX+b\in\mathbb C[X].
\]

Or
\[
X^2-aX+b=\left(X-\frac a2\right)^2-\frac{a^2-4b}{4}.
\]

En posant $a^2-4b=u+iv$ avec $u$ et $v$ réels on vérifie alors que :
\[
\frac12\left(\sqrt{u+\sqrt{u^2+v^2}}+i\sqrt{-u+\sqrt{u^2+v^2}}\right)^2=u+i|v|,
\]
(ne pas oublier que $\sqrt{v^2}=|v|$).

Donc si $v\ge0$,
\[
\delta=\frac1{\sqrt2}\left(\sqrt{u+\sqrt{u^2+v^2}}+i\sqrt{-u+\sqrt{u^2+v^2}}\right)
\]
\footnote{\textbf{校勘：}原式第二个根号内印作 $-u\sqrt{u^2+v^2}$。上一行恒等式给出的正确被开方数为 $-u+\sqrt{u^2+v^2}$；缺少加号时该式一般不能给出 $u+iv$ 的平方根。}
est dans $\mathbb C$ tel que $\delta^2=u+iv$, (et si $v<0$ on remplace $i$ par $-i$ dans $\delta$).

On a donc $\delta\in\mathbb C$ tel que
\[
X^2-aX+b=\left(X-\frac a2\right)^2-\delta^2
\]
d'où $z_{i_k}$ et $z_{j_k}$ dans $\mathbb C$ car du type
\[
\frac{a\pm\delta}{2}
\]
avec $a$ et $\delta$ dans $\mathbb C$.

\emph{Donc si $m=1$, $Q$ admet au moins 2 racines complexes, (éventuellement confondues).}

\emph{Si on suppose le résultat vrai pour $m-1$}, on repart comme précédemment, les $Q_k$ sont de degré $2^{m-1}\cdot($nombre impair$)$ donc ont des racines complexes. La même mise en forme aboutit à $Q$ admet 2 racines complexes, (distinctes ou non), et finalement on a bien justifié le fait que $P\in\mathbb C[X]$ admet au moins un zéro sur $\mathbb C$, (avec bien sûr $d^\circ P\ge1$).

Il est bien évident que, par application itérée de ce résultat et après factorisation des $(X-$ racine trouvée$)$, on a : tout polynôme $P$ de $\mathbb C[X]$ est \BellaCorrection{scindé}{scindée} sur $\mathbb C$. \bookend

\medskip\noindent\textsc{Corollaire 7.74.} --- \emph{Les seuls polynômes irréductibles de $\mathbb R[X]$ sont \BellaCorrectionNote{les polynômes de degré 1 et ceux de degré 2 sans racine réelle}{les constantes non nulles, les polynômes de degré 1 et ceux de degré 2 sans racine réelle}{D'après la définition 7.21, un élément irréductible doit être non inversible; or le Théorème 7.22 établit que les constantes réelles non nulles sont précisément des unités de $\mathbb R[X]$.}.}

Il est clair que les polynômes de degré 1 de $\mathbb R[X]$ sont irréductibles.

Soit $P$ de degré $\ge2$, irréductible sur $\mathbb R$. En particulier $P$ est sans racines \BellaCorrection{réelles}{réelle} (car si $\lambda$ racine, $P$ est divisible par $X-\lambda$, (Théorème 7.38)).

Mais $P$, dans $\mathbb C$, admet des racines. Soit $\lambda=u+iv$ avec $v\ne0$ un zéro de $P$. On a $P(\lambda)=0$ d'où en conjuguant $\overline{P(\lambda)}=0$ soit $P(\overline\lambda)=0$ puisque $P$ est à coefficients réels.

Comme dans $\mathbb C$ on a $\lambda\ne\overline\lambda$, les polynômes $X-\lambda$ et $X-\overline\lambda$ sont non associés (définition 7.32) donc $P$ est divisible, dans $\mathbb C[X]$, par
\[
(X-\lambda)(X-\overline\lambda)=(X-u-iv)(X-u+iv)
=(X-u)^2+v^2=X^2-2uX+(u^2+v^2)\in\mathbb R[X].
\]
Donc il existe $Q$ dans $\mathbb R[X]$ tel que
\[
P(X)=(X^2-2uX+u^2+v^2)Q(X),
\]
car, la division euclidienne, faite dans $\mathbb R[X]$ donne une identité
\[
P=(X^2-2uX+u^2+v^2)Q_1+R_1,
\]
valable aussi dans $\mathbb C[X]$ et la divisibilité de $P$ par $X^2-2uX+u^2+v^2$ dans $\mathbb C[X]$ donne $R_1=0$, d'où $Q_1=Q\in\mathbb R[X]$. Comme $P$ est irréductible, c'est que $d^\circ P=2$ et $Q$ est un scalaire réel non nul. \bookend

Enfin, pour conclure cet aspect constructif des grands ensembles mathématiques, passage de $\mathbb N$ à $\mathbb Z$ puis à $\mathbb Q$ par symétrisation de demi-groupe, puis de $\mathbb Q$ à $\mathbb R$ par complétion et de $\mathbb R$ à $\mathbb C$ par corps de rupture qui est en même temps une clôture algébrique, signalons que la clôture algébrique de $\mathbb Q$ n'est pas $\mathbb C$ : des nombres \BellaCorrection{transcendants}{trancendants} comme $e$ ou $\pi$ ne sont pas algébriques sur $\mathbb Q$. De plus, on démontre que le cardinal de la clôture algébrique de $\mathbb Q$ est celui de $\mathbb Q$ c'est-à-dire que ce corps est dénombrable.
